
\documentclass[10pt,a4paper]{article}
\usepackage{fontspec}
\setmainfont{Carlito}
\setsansfont{Carlito}
\usepackage[french]{babel}
\usepackage[a4paper,margin=18mm,headheight=17pt,footskip=12mm]{geometry}
\usepackage{microtype}
\usepackage{xcolor}
\definecolor{navy}{HTML}{153F56}
\definecolor{teal}{HTML}{157A86}
\definecolor{pale}{HTML}{EDF6F7}
\definecolor{warm}{HTML}{FFF4E4}
\definecolor{redsoft}{HTML}{FCECEC}
\definecolor{bodygray}{HTML}{445866}
\usepackage{hyperref}
\hypersetup{colorlinks=true,linkcolor=teal,urlcolor=teal,pdftitle={Analyse des SI - Révision et examen blanc}}
\usepackage{booktabs,tabularx,longtable,array,multicol}
\renewcommand{\arraystretch}{1.25}
\usepackage{enumitem}
\setlist[itemize]{leftmargin=*,itemsep=2pt,topsep=4pt}
\setlist[enumerate]{leftmargin=*,itemsep=2pt,topsep=4pt}
\usepackage{tikz}
\usetikzlibrary{positioning,arrows.meta,fit,calc,shapes.geometric,backgrounds}
\usepackage[most]{tcolorbox}
\tcbset{before skip=7pt,after skip=7pt,colback=pale,colframe=teal!40,arc=2mm,boxrule=.6pt,left=8pt,right=8pt,top=6pt,bottom=6pt}
\newtcolorbox{memo}[1][]{colback=pale,colframe=teal,title={À retenir},fonttitle=\bfseries,#1}
\newtcolorbox{warningbox}[1][]{colback=warm,colframe=orange!55!black,title={Piège à éviter},fonttitle=\bfseries,#1}
\newtcolorbox{exobox}[1][]{colback=white,colframe=navy!40,title={Exemple difficile},fonttitle=\bfseries,#1}
\newtcolorbox{answerbox}[1][]{colback=pale,colframe=teal!40,title={Justification / réponse},fonttitle=\bfseries,#1}
\usepackage{titlesec}
\titleformat{\section}{\Large\bfseries\color{navy}}{\thesection.}{.6em}{}
\titleformat{\subsection}{\large\bfseries\color{teal}}{\thesubsection}{.6em}{}
\titleformat{\subsubsection}{\normalsize\bfseries\color{navy}}{\thesubsubsection}{.6em}{}
\titlespacing*{\section}{0pt}{17pt}{8pt}
\titlespacing*{\subsection}{0pt}{12pt}{5pt}
\usepackage{fancyhdr,lastpage}
\pagestyle{fancy}
\fancyhf{}
\fancyhead[L]{\footnotesize\color{bodygray} M2 SIA -- Analyse des SI d'entreprise}
\fancyhead[R]{\footnotesize\color{bodygray} Révision UML}
\fancyfoot[L]{\footnotesize\color{bodygray} Support indépendant basé sur les cours 2026-2027}
\fancyfoot[R]{\footnotesize\color{bodygray} \thepage/\pageref{LastPage}}
\renewcommand{\headrulewidth}{.35pt}
\color{navy}
\setlength{\parindent}{0pt}
\setlength{\parskip}{5pt}
\newcommand{\tagword}[1]{\textcolor{teal}{\textbf{#1}}}
\newcommand{\good}{\textcolor{teal}{\textbf{OK}}}
\newcommand{\bad}{\textcolor{red!70!black}{\textbf{ERREUR}}}
\newcommand{\source}[1]{\par{\footnotesize\color{bodygray}\textit{Repère cours : #1}}}
\newcommand{\sep}{\smallskip\hrule\smallskip}
\tikzset{
 actor/.style={draw=navy,rounded corners=1pt,minimum width=16mm,minimum height=8mm,fill=white,font=\scriptsize,align=center},
 usecase/.style={draw=teal,ellipse,minimum width=27mm,minimum height=11mm,fill=pale,font=\scriptsize,align=center},
 cbox/.style={draw=teal,rounded corners=2pt,fill=pale,minimum width=26mm,minimum height=9mm,align=center,font=\scriptsize},
 dbox/.style={draw=teal,rectangle,fill=pale,minimum width=24mm,minimum height=8mm,align=center,font=\scriptsize},
 assoc/.style={-Stealth,draw=navy,thick},
 umlarrow/.style={-Stealth,draw=navy,dashed},
 life/.style={draw=bodygray,densely dashed},
 state/.style={draw=teal,rounded corners=4pt,fill=pale,minimum width=28mm,minimum height=11mm,align=center,font=\scriptsize}
}
\newcommand{\cover}[3]{
\thispagestyle{empty}
\vspace*{27mm}
{\color{teal}\rule{25mm}{3pt}}\par\vspace{12mm}
{\huge\bfseries\color{navy} #1\par}
\vspace{5mm}{\Large\color{teal}#2\par}
\vspace{13mm}\begin{tcolorbox}[colback=pale,colframe=teal!40,boxrule=.3pt]
#3
\end{tcolorbox}
\vfill
{\small\color{bodygray} Master 2 -- Systèmes d'information avancés\\
Analyse des SI d'entreprise -- iaelyon / Université Jean Moulin Lyon 3\\
Document d'entraînement original, non officiel -- octobre 2026}
\newpage
}

\begin{document}
\cover{Réviser toute l'analyse des SI}{Cours clair + cas difficiles + schémas UML}{\textbf{Objectif :} savoir lire une étude de cas longue, choisir les bons concepts, construire les diagrammes et expliquer ses choix en français simple.\\[3pt]
\textbf{Priorité pour l'écrit :} cas d'utilisation, séquence système N1, classes du domaine et paquetages, états, DCP, séquence N3 (MVC).\\[3pt]
\textbf{Compléments utiles :} description textuelle, activité, communication, démarche de projet, Scrum, vision, exigences, IA et SAFe.}
\tableofcontents
\newpage

\section{Ce que demande vraiment l'examen}
\begin{memo}
La séance 5 annonce un \textbf{examen écrit de 3 heures, 20 points, 5 questions, avec six fiches de diagrammes}. Les documents de cours sont autorisés. Le projet en groupe compte aussi pour 50\% de l'évaluation de l'UE. Le plan exact d'un futur sujet peut varier : ne mémorisez pas un seul cas.
\end{memo}
\begin{tabularx}{\linewidth}{@{}p{39mm}X@{}}
\toprule
Livrable central & Question à laquelle le diagramme répond \\
\midrule
Cas d'utilisation & Qui veut faire quoi avec le système ? \\
Séquence système N1 & Quels messages passent entre acteurs et système ? \\
Classes du domaine & Quels objets métier existent ? Quelles règles les lient ? \\
Paquetages & Comment ranger les classes sans dépendances circulaires ? \\
États & Comment une inscription, réservation, compte\ldots\ change d'état ? \\
DCP & Quelles classes de dialogue, contrôle et entité participent à un cas ? \\
Séquence N3 & Comment ces participants collaborent-ils dans le temps ? \\
\bottomrule
\end{tabularx}

\textbf{Autres notions du cours :} diagramme d'activité, diagramme de communication, modèle MVC, processus unifié, document de vision, risques, Scrum, exigences non fonctionnelles, usage critique de l'IA et architecture C4. Elles sont utiles dans les TD, le projet et les justifications. Le \textit{diagramme de communication} n'apparaît pas parmi les six fiches ciblées pour l'écrit, mais figure dans le TD « Organisation de courses de fond ».

\subsection{Méthode de gestion du temps (conseil d'entraînement)}
\begin{enumerate}
\item Lire \textbf{tout} le sujet et ses questions avant de dessiner (15 min).
\item Extraire les mots du métier, les rôles, les événements et les règles au brouillon (15 min).
\item Traiter les questions demandées ; commencer par la question 2 est le conseil explicite de la séance 5.
\item Écrire toute hypothèse avec un numéro : H1, H2\ldots
\item Garder les \textbf{15 dernières minutes} pour les multiplicités, branches d'erreur et cohérences entre diagrammes.
\end{enumerate}
\source{ASI\_C5, pages 2, 7--14 et 23.}

\section{Lire un très long énoncé sans se perdre}
\subsection{La méthode des cinq questions}
\begin{tabularx}{\linewidth}{@{}p{26mm}p{42mm}X@{}}
\toprule
Question & À extraire & Exemple (« centre de formation ») \\
\midrule
Qui ? & Acteurs, rôles externes & Élève, secrétaire, banque \\
Quoi ? & Objectifs métier & S'inscrire à une session \\
Quelles données ? & Noms du métier & Session, Inscription, Paiement \\
Quand ? & États, événements, délais & En attente, validée, after(48 h) \\
Quelle règle ? & Gardes, calculs, exceptions & Places limitées, certificat obligatoire \\
\bottomrule
\end{tabularx}
\begin{exobox}
« Un étudiant demande une place. La banque autorise la somme sans la débiter. Un responsable valide sous 48 h. Sinon, la demande expire. Un étudiant sur liste d'attente peut recevoir une nouvelle proposition. »\\
\textbf{Acteurs :} étudiant, responsable, banque (si externe).\\
\textbf{Classes candidates :} Étudiant, Session, Demande, AutorisationPaiement, Offre.\\
\textbf{États :} EnAttente, Confirmée, Refusée, Expirée ; des états de liste d'attente sont à justifier selon le périmètre.\\
\textbf{Règles :} pas de débit avant validation, délai de 48 h, gestion des places.
\end{exobox}
\begin{warningbox}
Ne copiez pas tous les noms en classes. « Bouton », « écran », « base de données », « email » (sans historique requis), « disponibilité » calculée : souvent des \textbf{faux candidats}. Le modèle du domaine parle du métier, pas de l'interface.
\end{warningbox}

\subsection{Construire une petite matrice de traçabilité}
\begin{tabularx}{\linewidth}{@{}p{27mm}p{31mm}X@{}}
\toprule
Exigence & Modèles concernés & Contrôle de cohérence \\
\midrule
Inscription avec paiement & UC, N1, classes, N3 & Même événement et même ordre d'autorisation. \\
Délai de 48 h & États, N1, règles & \texttt{after(48 h)}, annulation et résultat expliqués. \\
Certificat manquant & Texte, N1, N3, états & Branche d'erreur partout où elle compte. \\
Recherche d'une place & UC, DCP, N3 & Même nom de classe métier. \\
\bottomrule
\end{tabularx}

\section{Cas d'utilisation : objectifs et acteurs}
\subsection{Le bon raisonnement}
Un \textbf{acteur} est un rôle \emph{extérieur} au système étudié. Un acteur principal reçoit le service ; un acteur secondaire aide à le fournir (par exemple une banque ou un serveur externe). La même personne peut jouer plusieurs rôles. Un terminal ou une borne peut être un acteur si elle dialogue avec le système depuis l'extérieur de sa frontière.

Un \textbf{cas d'utilisation} est un objectif métier complet : « Réserver un atelier », « Terminer une réparation ». Il ne faut pas transformer « saisir le prénom » ou « cliquer sur OK » en cas autonome sans raison métier.

\begin{center}
\begin{tikzpicture}[node distance=20mm]
\node[actor] (client) {Client};
\node[usecase,right=21mm of client] (res) {Réserver\\un atelier};
\node[usecase,below=13mm of res] (ann) {Annuler\\une réservation};
\node[actor,right=25mm of res] (pay) {Banque\\externe};
\node[draw=teal,inner sep=8mm,fit=(res)(ann),label={[font=\scriptsize]above:Système de réservation}] (frame) {};
\draw[-,thick] (client) -- (res);\draw[-,thick] (client) -- (ann);\draw[-,thick] (res) -- (pay);
\end{tikzpicture}
\end{center}
\subsection{\texttt{include}, \texttt{extend}, généralisation}
\begin{center}
\begin{tikzpicture}[node distance=10mm and 27mm]
\node[usecase] (a) {Finaliser\\une commande};
\node[usecase,right=33mm of a] (b) {Vérifier\\le panier};
\draw[umlarrow] (a) -- node[above,font=\scriptsize]{\guillemotleft include\guillemotright} (b);
\node[usecase,below=21mm of a] (c) {Ajouter\\une garantie};
\node[usecase,right=33mm of c] (d) {Réserver\\un équipement};
\draw[umlarrow] (c) -- node[above,font=\scriptsize]{\guillemotleft extend\guillemotright} (d);
\end{tikzpicture}
\end{center}
\textbf{Include :} le cas de base appelle obligatoirement un comportement réutilisable ; \textbf{la flèche pointe vers le cas inclus}.\\
\textbf{Extend :} comportement complémentaire activé sous condition ; \textbf{la flèche pointe vers le cas de base}.\\
\textbf{Généralisation :} le rôle spécialisé hérite des interactions du rôle général ; triangle blanc vers le rôle général.
\begin{warningbox}
Une étape obligatoire n'est pas automatiquement un \texttt{include}. Le professeur sanctionne les \texttt{include} créés seulement pour découper chaque clic. La règle « paiement toujours réalisé » dépend du scénario : \emph{autorisation} et \emph{débit} ne sont pas la même chose.
\end{warningbox}
\begin{exobox}
Un gérant est aussi mécanicien et peut acheter comme client. On peut représenter \texttt{Gérant} et \texttt{Mécanicien} comme des \emph{rôles}, avec leurs cas propres. On ne crée pas un cas \emph{Créer un gérant} sauf objectif métier réel. Une machine qui envoie ses mesures peut être acteur secondaire si elle est hors frontière du système.
\end{exobox}
\source{ASI\_C1, partie cas d'utilisation ; ASI\_C5, page 8 ; TD Station multi-énergies et Agence de voyages.}

\section{Description textuelle détaillée : ne jamais l'oublier}
\begin{warningbox}
Le \textbf{travail 2} du TD \emph{Organisation de courses de fond} demande cette description. Même lorsque l'écrit note surtout des diagrammes, la description donne les messages de la N1. \textbf{Une liste de clics ou trois phrases ne suffit pas}. Le cours ASI\_C1, page 36, montre un exemple complet de retrait d'argent.
\end{warningbox}
\subsection{Les rubriques exactes et leur rôle}
\begin{longtable}{@{}p{37mm}p{102mm}@{}}
\toprule
Rubrique & Ce qu'il faut écrire, en français très simple \\
\midrule\endhead
\textbf{Identification} & Nom du cas, \textbf{but}, acteur principal, acteurs secondaires ; date, version, responsable si demandés.\\
\textbf{Déclencheur} & Événement \emph{qui commence} le cas : « Le concurrent demande une inscription sur internet. »\\
\textbf{Préconditions} & Ce qui doit déjà être vrai \emph{avant de commencer}. Exemple : l'épreuve a été créée et publiée dans le système.\\
\textbf{Enchaînement nominal} & Étapes numérotées acteur/système, jusqu'au succès. Dire \emph{qui fait quoi} à chaque étape.\\
\textbf{Enchaînements alternatifs} & Autre chemin possible, avec une référence à l'étape où il démarre : paiement par chèque, profil déjà connu, inscription par club.\\
\textbf{Exceptions} & Échec et réponse du système : quota atteint, paiement refusé, pièce non conforme. Écrire si on reprend le scénario ou s'il s'arrête.\\
\textbf{Postconditions} & Ce qui est vrai \emph{après le cas} ; distinguer \textbf{succès}, \textbf{échec} et parfois \textbf{dossier en cours}.\\
\textbf{Rubriques optionnelles} & Règles métier (tarif, âge, délai), exigences non fonctionnelles (confidentialité, fiabilité), IHM et hypothèses notées H1, H2.\\
\bottomrule
\end{longtable}
\begin{memo}
\textbf{Précondition vs garde vs postcondition :}\newline
Précondition : « L'épreuve existe avant le début du cas. »\newline
Garde/contrôle pendant le scénario : « [places disponibles ET date au plus tard J--15] ».\newline
Postcondition : « L'inscription est DEFINITIVE » \emph{ou} « la demande est ANNULEE ».\newline
\textbf{Piège :} le quota et la limite de date sont des règles à \emph{tester}, pas seulement des préconditions cachées : sinon on oublie la branche de refus.
\end{memo}
\subsection{Exemple complet corrigé : « S'inscrire à une course sur internet »}
\begin{exobox}
\textbf{Identification}\newline
\textbf{Nom :} S'inscrire sur internet. \textbf{But :} obtenir une inscription définitive à une course.\newline
\textbf{Principal :} Concurrent, pour une course individuelle. \textbf{Secondaires :} Banque (paiement par carte) et Secrétariat (chèque et certificat).\newline
\textbf{Déclencheur :} Le concurrent demande à participer à une course depuis le site.\newline
\textbf{Préconditions :} La course existe et le système d'inscription est accessible. \emph{Le nombre de places et la date limite seront vérifiés dans le déroulement.}
\end{exobox}
\textbf{Enchaînement nominal (course individuelle, paiement par carte) :}
\begin{enumerate}[label=\textbf{\arabic*.}]
\item Le concurrent choisit la course. Le système affiche les conditions et le tarif.
\item Le système vérifie \textbf{le quota}, \textbf{la date} (au plus tard J--15) et \textbf{l'âge} lié à l'épreuve.
\item Le concurrent indique son identité et ses coordonnées. Le système retrouve les informations connues, si elles existent.
\item Le système enregistre la \texttt{PREINSCRIPTION}.
\item Le concurrent choisit de payer par carte. Le système transmet le paiement à la banque.
\item La banque donne son accord. Le système valide l'inscription financière : statut \texttt{FINANCIERE}.
\item Le concurrent envoie un certificat médical de moins d'un an \textbf{par courrier}. Le secrétariat reçoit et contrôle la pièce.
\item Le secrétariat valide la conformité. Le système place le dossier en \texttt{DEFINITIVE} et informe le concurrent.
\end{enumerate}
\textbf{Alternatives (avec point de retour) :}
\begin{itemize}
\item \textbf{A1, étape 3 : profil déjà connu.} Le système réutilise les données, le concurrent les vérifie, puis retour à l'étape 4. \textbf{La pièce médicale reste à contrôler.}
\item \textbf{A2, étapes 5--6 : paiement par chèque.} Le concurrent envoie son chèque par courrier. Le secrétariat vérifie qu'il est conforme. Si oui, retour à l'étape 7 ; le statut reste \texttt{PREINSCRIPTION} tant que le paiement n'est pas validé.
\item \textbf{A3, étapes 3--8 : triathlon / vélo.} \textbf{Le Club} inscrit son membre. Il doit être un club de la discipline. \textbf{La licence remplace le certificat médical.} Aucune inscription sur place pour ces disciplines.
\end{itemize}
\textbf{Exceptions (avec effet sur les données) :}
\begin{itemize}
\item \textbf{E1, étape 2 :} \texttt{[quota atteint OU date après J--15 OU non admissible]}. La demande est refusée avant création de préinscription ; le système explique pourquoi.
\item \textbf{E2, étape 6 :} \texttt{[paiement non valide]}. L'inscription est \texttt{ANNULEE} et le concurrent est prévenu, comme le dit le sujet.
\item \textbf{E3, étape 8 :} \texttt{[certificat non conforme]}. L'inscription est \texttt{ANNULEE} et le concurrent est prévenu.
\end{itemize}
\textbf{Postconditions :}
\begin{itemize}
\item \textbf{Succès :} une \texttt{Inscription} existe au statut \texttt{DEFINITIVE}, rattachée à une \texttt{Epreuve} et un \texttt{Concurrent} ; le paiement et le justificatif sont acceptés.
\item \textbf{Intermédiaire :} le dossier peut être en attente de chèque ou de certificat. Il reste au stade de préinscription ou d'inscription financière.
\item \textbf{Échec :} aucune préinscription si E1 ; inscription \texttt{ANNULEE} et information envoyée si E2 ou E3.
\end{itemize}
\textbf{Règles métier :} 16 ans minimum pour une course libre, 18 pour le semi-marathon, 20 pour le marathon ; club spécialisé pour triathlon/vélo ; quota non dépassé ; limite J--15 ; certificat de moins d'un an pour les courses individuelles. \textbf{Exigence non fonctionnelle (proposition, non imposée mot à mot) :} protéger les données personnelles et médicales.
\begin{warningbox}
\textbf{Ne pas inventer dans ce cas :} un \textit{after(24 h)}, un débit différé après validation du gestionnaire, ou un téléchargement obligatoire du certificat sur le site. Ces éléments apparaissent dans \emph{d'autres} études de cas, \textbf{pas} dans le parcours internet du TD courses de fond.
\end{warningbox}
\subsection{Les trois voies : surtout ne pas les mélanger}
\begin{tabularx}{\linewidth}{@{}p{28mm}XXX@{}}
\toprule
Critère & Internet & Sur place & Courrier \\
\midrule
Qui saisit ? & Concurrent ; club pour triathlon/vélo & Secrétariat (H1) & Secrétariat (H1), à la réception\\
Date clé & Préinscription au plus tard J--15 & Terminer \textbf{1 h avant départ} & Terminer au plus tard J--5\\
Règlement & Carte ou chèque & Espèces, carte ou chèque & Chèque / règlement envoyé\\
Pièce & Certificat envoyé par courrier ; licence si club & Certificat contrôlé à nouveau & Certificat joint ; licence si club\\
Pièce manquante & Le dossier attend la suite des étapes si possible & Les 3 étapes doivent être validées & Reste à PREINSCRIPTION ou FINANCIERE ; lettre de relance\\
Triathlon/vélo & Par club & \textbf{Interdit} & Par club\\
\bottomrule
\end{tabularx}
\textbf{Passage au diagramme N1 :} chaque étape est un message, chaque A1--A3 ou E1--E3 devient une branche \texttt{alt}/\texttt{opt}/\texttt{break} lorsque c'est pertinent. En N3, la vue transmet au contrôle et les entités effectuent les règles métier.\newline
\source{ASI\_C1, page 36 (modèle du retrait d'argent) ; Organisation de courses de fond, pages 3--4 ; ASI\_C5, pages 4--5 et 9.}

\section{Séquence système N1 : un système en boîte noire}
\subsection{Ce qu'il faut dessiner}
Des lignes de vie pour les acteurs extérieurs et \textbf{une seule pour le système étudié}. Les \emph{messages} vont dans l'ordre chronologique. Les réponses importantes apparaissent en pointillés. Aucun écran, contrôleur, DAO ou base interne.

\begin{center}
\begin{tikzpicture}[x=1cm,y=.72cm]
\foreach \x/\txt in {0/Client,4/Système,8/Banque} {
 \node[dbox,minimum width=23mm] at (\x,0) {\txt};
 \draw[life] (\x,-.45)--(\x,-6.5);
}
\draw[assoc] (0,-1) -- node[above,font=\scriptsize]{1: demanderRéservation()} (4,-1);
\draw[assoc] (4,-2) -- node[above,font=\scriptsize]{2: autoriser(montant)} (8,-2);
\draw[<- ,dashed] (4,-3) -- node[above,font=\scriptsize]{3: accepté/refusé} (8,-3);
\draw[teal,thick] (3.4,-3.5) rectangle (8.6,-5.5);
\node[font=\scriptsize,anchor=north west] at (3.5,-3.5) {alt [autorisation OK] / [refusée]};
\draw[<- ,dashed] (0,-4.2) -- node[above,font=\scriptsize]{4a: demande en attente} (4,-4.2);
\draw[<- ,dashed] (0,-5) -- node[above,font=\scriptsize]{4b: refus du paiement} (4,-5);
\end{tikzpicture}
\end{center}
\subsection{Fragments combinés : choisir correctement}
\begin{tabularx}{\linewidth}{@{}p{18mm}p{35mm}X@{}}
\toprule
Fragment & Utilité & Exemple \\
\midrule
\texttt{alt} & Plusieurs issues exclusives & [paiement OK] / [paiement refusé] \\
\texttt{opt} & Issue conditionnelle unique & [client demande un reçu] \\
\texttt{loop} & Répétition & [pour chaque exemplaire], scanner \\
\texttt{break} & Fin anticipée du scénario & [carte inconnue], stopper l'emprunt \\
\texttt{par} & Branches en parallèle & Deux vérifications indépendantes (si justifié) \\
\bottomrule
\end{tabularx}
\begin{exobox}
La médiathèque vérifie chaque objet apporté. Un objet dépassant le quota est ignoré mais le suivant est testé. Il faut \texttt{loop} sur les objets et, \emph{à l'intérieur}, \texttt{alt} [quota OK] / [quota atteint]. Ce n'est pas un \texttt{break} global : le cas continue avec l'objet suivant.
\end{exobox}
\source{ASI\_C1, section séquences ; ASI\_C5, page 9 ; Exercices\_03, emprunts du tiers-lieu culturel.}

\section{Diagramme d'activité : flux de travail et partitions}
Un diagramme d'activité montre les \textbf{actions et décisions} d'un processus. Il est utile pour une inscription en trois étapes ou un prêt de plusieurs objets.
\begin{center}
\begin{tikzpicture}[node distance=8mm and 15mm]
\node[circle,fill=navy,inner sep=3pt] (start) {};
\node[cbox,below=of start] (s) {Recevoir la demande};
\node[diamond,draw=teal,aspect=2.5,below=of s,inner sep=1pt,fill=pale,font=\scriptsize] (d) {Dossier valide ?};
\node[cbox,below left=13mm and 19mm of d] (bad) {Demander les pièces};
\node[cbox,below right=13mm and 19mm of d] (good) {Enregistrer};
\node[circle,draw=navy,thick,minimum size=7mm,below=28mm of d] (end) {};
\node[circle,fill=navy,inner sep=2pt] at (end.center) {};
\draw[assoc] (start)--(s);\draw[assoc] (s)--(d);
\draw[assoc] (d)--node[above left,font=\scriptsize]{[non]}(bad);
\draw[assoc] (d)--node[above right,font=\scriptsize]{[oui]}(good);
\draw[assoc] (bad)|-(end);\draw[assoc] (good)|-(end);
\end{tikzpicture}
\end{center}
\textbf{Initial/final :} point noir / point noir cerclé. \textbf{Action :} rectangle arrondi. \textbf{Décision :} losange avec gardes [ ]. \textbf{Fork/Join :} barres épaisses pour créer puis synchroniser deux branches. \textbf{Partitions (swimlanes) :} séparer les actions du client, du secrétariat, de la banque.
\begin{warningbox}
\emph{Décision} (un seul chemin choisi) \(\neq\) \emph{fork} (plusieurs chemins en parallèle). Un join attend normalement les branches du fork ; n'en ajoutez pas pour un simple choix oui/non.
\end{warningbox}
\source{ASI\_C1, diagrammes d'activité et partitions ; cas Cloud gaming et projet.}

\section{Classes du domaine : le chapitre qui piège le plus}
\subsection{Classe ou attribut ? Les six tests du cours}
Pour chaque nom extrait du texte, demandez-vous : (1) plusieurs exemplaires ? (2) données propres ? (3) cycle de vie propre ? (4) ensemble fermé de valeurs ? (5) information portée par le \emph{lien} ? (6) identité propre du lien ? Ces tests orientent vers classe, attribut, énumération, classe d'association ou classe ordinaire.

\begin{tabularx}{\linewidth}{@{}p{38mm}X@{}}
\toprule
Expression du texte & Choix raisonnable / justification \\
\midrule
Date de début & Attribut de Réservation : valeur simple. \\
Statut = {attente, OK, refus} & Énumération ; mêmes valeurs que dans les états. \\
Tarif par catégorie et durée & Classe Tarif : plusieurs attributs, évolue dans le métier. \\
Date d'une inscription d'un élève à un cours & Attribut du lien si une seule inscription par couple ; classe autonome si plusieurs inscriptions, identité ou cycle de vie. \\
Disponibilité & Calcul métier ; souvent \texttt{/disponible}, pas une classe. \\
ÉcranRecherche / BaseDeDonnées & Hors du modèle métier ; elles reviennent en conception. \\
\bottomrule
\end{tabularx}
\subsection{Attributs, opérations et visibilité}
\textbf{Analyse :} nom métier, attributs utiles, opérations métier (\texttt{calculerPrix()}, \texttt{verifierEligibilite()}). \textbf{Conception :} types techniques, visibilité \texttt{+/-}, constructeur, accesseurs, dépendances techniques. Ne mettez pas par réflexe un \texttt{getNom()} ou une classe \texttt{DAO} dans le diagramme \emph{du domaine}.

\subsection{Multiplicités à lire par l'autre extrémité}
\begin{center}
\begin{tikzpicture}
\node[dbox] (ad) {Adhérent};\node[dbox,right=48mm of ad] (re) {Réservation};
\draw[-,thick] (ad)--node[above,pos=.12,font=\scriptsize]{1}node[above,pos=.88,font=\scriptsize]{0..*}(re);
\end{tikzpicture}
\end{center}
\textbf{Lecture :} pour \emph{une réservation}, \texttt{1} adhérent ; pour \emph{un adhérent}, \texttt{0..*} réservations.\\
\texttt{0..1} facultatif unique ; \texttt{1..*} au moins une occurrence ; \texttt{2..5} borné ; \texttt{*} équivaut à \texttt{0..*}.

\subsection{Association, héritage, agrégation, composition}
\begin{itemize}
\item \textbf{Association :} deux concepts sont liés (un étudiant suit un cours).
\item \textbf{Généralisation / héritage :} « une Voiture est un Véhicule ». Triangle blanc vers la classe mère. Utiliser seulement si la substitution a du sens.
\item \textbf{Composition :} losange noir côté \emph{tout}. Une ligne de facture n'a pas de vie autonome hors de la facture ; forte dépendance de cycle de vie.
\item \textbf{Agrégation partagée :} losange blanc ; le composant peut exister ou être partagé indépendamment. À employer avec prudence : une simple association peut suffire.
\item \textbf{Classe d'association :} le lien porte des attributs (ex. date d'affectation à une équipe). Si une même paire peut avoir plusieurs liens datés, préférer souvent une classe ordinaire identifiable.
\end{itemize}
\begin{center}
\begin{tikzpicture}[node distance=19mm]
\node[dbox] (f) {Facture};\node[dbox,right=43mm of f] (l) {LigneFacture};
\draw[-,thick] (f) -- node[above,pos=.9,font=\scriptsize]{1..*} (l);
\node[diamond,draw=navy,fill=navy,inner sep=2.4pt,rotate=90] at ($(f.east)+(3mm,0)$) {};
\end{tikzpicture}
\end{center}
\begin{warningbox}
Une classe « prix1, prix2, prix3 » cache une classe métier manquante. La composition ne se choisit pas avec l'intuition « contient ». Testez \textbf{destruction, partage, création}. Et renseignez \emph{les deux extrémités} de chaque association.
\end{warningbox}
\source{ASI\_C2, pages 5--30 ; ASI\_C5, page 11 ; TD application hôtelière et banque.}

\section{Paquetages : découpage propre et dépendances}
Un \textbf{paquetage} regroupe des concepts du même domaine de responsabilité. Un paquetage A dépend de B si A \emph{utilise} des classes de B. Le but est un modèle cohérent, avec peu de liens et \textbf{sans cycle de dépendances}.
\begin{center}
\begin{tikzpicture}[node distance=11mm and 19mm]
\node[dbox,minimum width=29mm] (cli) {Clients};
\node[dbox,right=of cli,minimum width=35mm] (res) {Réservations};
\node[dbox,right=of res,minimum width=29mm] (pai) {Paiements};
\draw[umlarrow] (res)--(cli); \draw[umlarrow] (pai)--(res);
\end{tikzpicture}
\end{center}
Les flèches vont \emph{du paquetage utilisateur vers celui qui est utilisé}. L'exemple montre une orientation possible, pas une règle universelle sur les paiements.
\begin{exobox}
Un modèle place \texttt{Client} dans \texttt{FrontOffice}, \texttt{Paiement} dans \texttt{BackOffice}, \texttt{Réservation} dans \texttt{Mobile}. C'est fragile : un même concept métier serait dispersé selon les écrans. Meilleur principe : \emph{comptes}, \emph{réservations}, \emph{facturation}.\end{exobox}

\section{Diagramme d'états : vrai cycle de vie}
Un état décrit la situation \emph{durable} d'un objet (\texttt{EnAttente}, \texttt{Confirmée}). Une transition écrit \textbf{événement [garde] / action}. \texttt{after(durée)} représente l'écoulement du temps. Les états doivent être repris \textbf{mot pour mot} dans une énumération du domaine.
\begin{center}
\begin{tikzpicture}[node distance=14mm and 10mm]
\node[circle,fill=navy,inner sep=2.5pt] (i) {};
\node[state,right=of i] (a) {Préinscrite};
\node[state,right=27mm of a] (p) {PaiementValidé};
\node[state,right=26mm of p] (c) {Confirmée};
\node[state,below=23mm of p] (x) {Annulée};
\draw[assoc] (i)--(a);
\draw[assoc] (a)--node[above,font=\scriptsize]{paiement accepté}(p);
\draw[assoc] (p)--node[above,font=\scriptsize]{certificat conforme}(c);
\draw[assoc] (p)--node[right,font=\scriptsize,align=left]{certificat refusé\\/ notifier()}(x);
\draw[assoc] (a) to[bend right=20] node[below left,font=\scriptsize]{paiement refusé} (x);
\end{tikzpicture}
\end{center}
\textbf{Attention, cas réel « Organisation de courses de fond » :} la préinscription ne peut commencer que si la capacité n'est pas atteinte, au plus tard 15 jours avant la course ; le paiement est contrôlé ; un certificat de moins d'un an valide l'inscription définitive. Pour le courrier, le dossier peut \textbf{rester} en préinscription ou inscription financière si une pièce manque. Pour le triathlon et le vélo, la licence remplace le certificat et l'inscription passe par le club. Un vrai schéma doit représenter ces retours et branches, sans inventer un parcours unique pour tout le monde.
\begin{warningbox}
Ne confondez pas \emph{événement} (« recevoir le paiement »), \emph{garde} (« [paiement valide] »), \emph{action} (« enregistrer le paiement ») et \emph{état} (« PaiementValidé »). Les flèches sans libellé font perdre des points. Ne supprimez pas les transitions de retour imposées par le texte.
\end{warningbox}
\source{ASI\_C2, partie états ; ASI\_C5, page 12 ; Organisation de courses de fond, pages 3--6.}

\section{DCP : classes participantes au cas d'utilisation}
Le \textbf{DCP} contient trois familles de classes, en restant lié à \emph{un cas} :
\begin{itemize}
\item \textbf{Dialogue / Boundary :} une classe par écran (recherche, saisie, confirmation).
\item \textbf{Contrôle :} une classe de contrôle par cas d'utilisation ; elle organise les échanges.
\item \textbf{Entité :} les classes métier utilisées, avec exactement les noms du diagramme de classes du domaine.
\end{itemize}
\begin{center}
\begin{tikzpicture}[node distance=11mm and 14mm]
\node[actor] (ac) {Adhérent};
\node[dbox,right=of ac] (v) {«dialogue»\\ÉcranRéservation};
\node[dbox,right=of v] (co) {«contrôle»\\ContrôleRéserver};
\node[dbox,right=of co] (en) {«entité»\\Réservation};
\draw[-,thick] (ac)--(v);\draw[-,thick] (v)--(co);\draw[-,thick] (co)--(en);
\end{tikzpicture}
\end{center}
\textbf{Règle absolue du cours :} pas de lien direct \texttt{ÉcranRéservation -- Réservation}. L'acteur dialogue uniquement avec l'interface. \emph{Le DCP n'est pas encore un diagramme chronologique} : il montre les collaborations et attributs/opérations pertinents.
\begin{exobox}
Un projet a trois écrans : recherche, choix, confirmation. Un DCP avec \emph{un seul écran générique} peut être incomplet si les trois maquettes figurent dans l'énoncé. Trois dialogues, un contrôle du cas, entités du domaine ; ne créez pas trois « ContrôleurGlobaux » sans objectif.
\end{exobox}
\source{ASI\_C3, DCP ; ASI\_C5, page 13.}

\section{Séquence N3 (MVC) : ouvrir la boîte noire}
La N3 répond au \emph{même cas} que la N1, mais montre \textbf{vue, contrôle, modèle}. Le modèle possède les règles du métier ; le contrôle coordonne les objets ; la vue affiche et reçoit les actions.
\begin{center}
\begin{tikzpicture}[x=1cm,y=.68cm]
\foreach \x/\t in {0/Client,2.6/Vue,5.2/Contrôle,7.8/Réservation} {
 \node[dbox,minimum width=22mm] at (\x,0) {\t};\draw[life] (\x,-.43)--(\x,-7);
}
\draw[assoc] (0,-1)--node[above,font=\scriptsize]{1: demander()}(2.6,-1);
\draw[assoc] (2.6,-2)--node[above,font=\scriptsize]{2: soumettre()}(5.2,-2);
\draw[assoc] (5.2,-3)--node[above,font=\scriptsize]{3: vérifier()}(7.8,-3);
\draw[<-,dashed] (5.2,-4)--node[above,font=\scriptsize]{4: résultat}(7.8,-4);
\draw[teal] (4.4,-4.5) rectangle (8.6,-6.15);
\node[font=\scriptsize,anchor=north west] at (4.5,-4.6) {alt [valide] / [refusé]};
\draw[<-,dashed] (2.6,-5.5)--node[above,font=\scriptsize]{5: état ou erreur}(5.2,-5.5);
\draw[<-,dashed] (0,-6.6)--node[above,font=\scriptsize]{6: afficher()}(2.6,-6.6);
\end{tikzpicture}
\end{center}
\begin{warningbox}
Trois violations MVC à connaître : \textbf{Vue \(\rightarrow\) Entité} sans contrôle ; contrôle qui calcule seul les règles du métier ; entité qui commande l'affichage. La N3 doit reprendre les échanges visibles de la N1 \textbf{et la branche d'erreur}.
\end{warningbox}
\textbf{Niveaux :} N1 = système noir ; N2 = objets du domaine visibles (vue d'analyse intermédiaire) ; N3 = classes participantes et conception MVC, selon la progression du cours.
\source{ASI\_C3, partie MVC et séquences ; ASI\_C5, page 10.}

\section{Diagramme de communication : messages numérotés}
Même idée de collaboration qu'une séquence N3, mais sans axe du temps vertical. On pose les participants, on trace les liens, on inscrit \textbf{1:, 2:, 3:\ldots} sur les messages, avec flèche de direction. Les appels imbriqués peuvent utiliser \texttt{1.1}, \texttt{1.2}.
\begin{center}
\begin{tikzpicture}[node distance=17mm and 28mm]
\node[dbox] (v) {ÉcranRecherche};\node[dbox,right=of v] (c) {ContrôleRecherche};
\node[dbox,below=of c] (r) {Classement};
\draw[-,thick] (v)--node[above,font=\scriptsize]{1: rechercher()}(c);
\draw[-,thick] (c)--node[right,font=\scriptsize]{2: trouver()}(r);
\draw[assoc] ($(v.east)+(1mm,-3mm)$)--($(c.west)+(-1mm,-3mm)$);
\draw[assoc] ($(c.south)+(-3mm,-1mm)$)--($(r.north)+(-3mm,1mm)$);
\end{tikzpicture}
\end{center}
Le TD \emph{Organisation de courses de fond} demande ce modèle pour \emph{rechercher un classement}. C'est une extension utile de la N3 : même cas, mêmes entités et mêmes opérations, \textbf{messages ordonnés}.

\section{Processus unifié, vision, exigences et risques}
\subsection{Processus unifié en quatre phases}
\begin{tabularx}{\linewidth}{@{}p{29mm}X@{}}
\toprule
Phase & But en une phrase \\
\midrule
Inception & Pourquoi le projet ? Périmètre, utilisateurs, vision, risques, coût. \\
Élaboration & Réduire les risques critiques et définir une architecture qui tient. \\
Construction & Construire le produit par itérations, tester et intégrer. \\
Transition & Déployer, former, corriger et faire adopter le système. \\
\bottomrule
\end{tabularx}
La démarche est \textbf{itérative} : le diagramme d'états peut faire corriger les classes ; les classes peuvent faire corriger les cas d'utilisation.

\subsection{Document de vision : quoi remplir ?}
Le gabarit du cours comprend notamment : \textbf{position du problème}, position du produit, parties prenantes/acteurs, périmètre/hors périmètre, exigences non fonctionnelles mesurables, contraintes/hypothèses, risques principaux, itérations/jalons, maquettes. Il explique \emph{pourquoi} et \emph{pour qui} ; il ne remplace pas les descriptions détaillées des cas.

\begin{exobox}
\bad{} « L'application doit être rapide et très sécurisée. »\\
\good{} « 95\% des recherches répondent en moins de 2 s ; les paiements passent par un prestataire certifié ; tout changement de rôle est tracé. » Ce sont des exigences \emph{vérifiables} ; elles influencent le modèle et l'architecture.
\end{exobox}
\subsection{Contrainte, hypothèse et risque : ne pas confondre}
\textbf{Contrainte :} imposée et connue (loi, délai, budget). \textbf{Hypothèse :} supposée vraie mais non confirmée (H1 : un compte ne peut avoir qu'une réservation active par créneau). \textbf{Risque :} événement possible qui peut gêner le projet (indisponibilité de l'API bancaire). Une fiche de risque précise probabilité, gravité, parade, responsable et moment de levée.
\source{ASI\_C4 et Document de vision.pdf.}

\section{Scrum, MVP, story map, SAFe : théorie à connaître}
\subsection{Scrum : 3 responsabilités, 5 événements, 3 artefacts}
\begin{tabularx}{\linewidth}{@{}p{33mm}X@{}}
\toprule
Élément & Sens simple \\
\midrule
Product Owner & Responsable de la valeur et de l'ordre du Product Backlog. \\
Scrum Master & Aide à comprendre et appliquer Scrum, enlève les obstacles. \\
Developers & Créent l'incrément utilisable ; ils s'organisent eux-mêmes. \\
Sprint & Cycle court contenant les autres événements. \\
Sprint Planning & Choisir pourquoi, quoi et comment travailler pendant le sprint. \\
Daily Scrum & 15 minutes pour inspecter l'avancement vers le Sprint Goal. \\
Sprint Review & Examiner le résultat et adapter la suite avec les parties prenantes. \\
Sprint Retrospective & Améliorer le fonctionnement de l'équipe. \\
Product Backlog / Product Goal & Travail à venir ordonné / objectif du produit. \\
Sprint Backlog / Sprint Goal & Plan du sprint / objectif du sprint. \\
Increment / Definition of Done & Résultat utilisable / qualité nécessaire pour être terminé. \\
\bottomrule
\end{tabularx}
\textbf{Story map :} ranger les activités du parcours utilisateur puis les découper en tâches / versions. \textbf{MVP :} plus petite version donnant une valeur vérifiable et permettant d'apprendre. Ce n'est pas « la version mal finie ».
\begin{warningbox}
Scrum n'est pas un cycle en V décoré d'un Daily Scrum. Un Product Owner sans droit de décider et un backlog figé un an ne suffisent pas à rendre un projet agile. \emph{Affinage du backlog} = activité continue, pas sixième événement officiel.
\end{warningbox}
\textbf{SAFe :} un cadre de Lean-Agile à grande échelle, pour aligner stratégie, équipes et livraisons. Distinguer un besoin d'échelle réel d'une organisation qui ajoute seulement des réunions. Le livre SAFe joint est un complément, pas une fiche UML parmi les six de l'examen.
\source{ASI\_C4 ; SAFe\_Explained\_Ebook\_FR.pdf (complément).}

\section{MVC, architecture et décisions de conception}
\textbf{MVC :} Modèle = données et règles du métier ; Vue = présentation ; Contrôleur = coordination des actions. En conception on peut avoir d'autres composants techniques, mais ils ne doivent pas polluer le diagramme \emph{du domaine}.\\
\textbf{C4 (vue architecture) :} contexte d'un système, conteneurs applicatifs, composants et code ; ces niveaux servent à discuter du système à différentes échelles. Ils ne remplacent ni le diagramme de classes ni la séquence N3.\\
\textbf{Non-fonctionnel et RGPD :} volumétrie, disponibilité, confidentialité, conservation des données, fonctionnement hors ligne ont des conséquences concrètes : archivage, délai, accès, suppression/anonymisation, synchronisation.
\begin{exobox}
Une application de location transmet la position des voitures toutes les deux minutes. Le texte impose \emph{uniquement la dernière position en location normale}, et l'historique \emph{seulement} après dépassement de l'heure de restitution ou vol déclaré. Il serait faux de conserver toujours un \texttt{TrajetClient} « pour faire de l'IA ». Le respect du périmètre et de la vie privée modifie le modèle.
\end{exobox}

\section{IA, qualité des modèles et traçabilité}
Le cours permet l'usage d'IA \textbf{dans le projet s'il est déclaré} ; l'étudiant est responsable de la qualité du rendu et doit justifier ses choix. Pour l'examen, suivez strictement les consignes de surveillance et les ressources effectivement autorisées ; l'autorisation d'apporter des documents n'est pas une permission générale d'utiliser des outils numériques.

\subsection{Sept liens à vérifier selon la séance 5}
\begin{enumerate}
\item Chaque cas décrit a sa N1 ; les alternatives deviennent des fragments.
\item Chaque maquette sert un cas ; chaque écran devient dialogue dans un DCP.
\item Chaque classe du domaine sert vraiment à quelque chose.
\item Chaque état = une valeur de l'énumération d'état, même orthographe.
\item Chaque entité du DCP garde le nom du domaine.
\item La N3 conserve les échanges du cas visibles en N1.
\item Les hypothèses H1, H2\ldots\ sont écrites et liées aux modèles concernés.
\end{enumerate}
\begin{warningbox}
Les erreurs qui coûtent cher : multiplicités absentes, CRUD comme cas d'utilisation, classes techniques dans le domaine, écran relié directement à une entité, cycle d'états sans retour/délai, contrôleur global, flèches \texttt{include/extend} inversées et \textbf{incohérence entre deux diagrammes}.
\end{warningbox}
\source{ASI\_C5, pages 5 et 14.}

\section{Trois mini-cas difficiles corrigés}
\subsection{Cas A : voyage, facture à la demande}
\textbf{Énoncé :} Une agence propose des transports inclus dans un séjour. Après la réservation, le client peut demander une facture détaillée. Les disponibilités viennent d'une API externe et ne sont pas conservées.\\
\textbf{Solution raisonnée :} \texttt{Réserver un séjour} est un cas métier. Le transport n'est pas un cas \emph{autonome} si le texte interdit de le réserver seul. La facture détaillée est \texttt{extend} sous condition, flèche vers le cas de base choisi. L'API partenaire est acteur secondaire. Ne créez pas de classe \texttt{Disponibilité} persistante puisqu'elle est interrogée à la demande.

\subsection{Cas B : hôtel, réservation multi-chambres}
\textbf{Énoncé :} Le client réserve plusieurs chambres, entre deux dates. Une chambre peut être louée à beaucoup de dates, jamais à deux clients sur des périodes qui se chevauchent. Une plateforme transmet ses ventes seulement par lots.\\
\textbf{Solution raisonnée :} \texttt{Réservation} porte dates/statut et lie client aux chambres. Pour plusieurs chambres par réservation, justifier une association à plusieurs chambres ou une \texttt{LigneRéservation} si des données propres à chaque ligne sont nécessaires. Le non-chevauchement est une \emph{règle temporelle}, pas une simple multiplicité \texttt{0..1}. La synchronisation par lots est une contrainte fonctionnelle/non fonctionnelle à signaler, sans inventer de disponibilité instantanée.

\subsection{Cas C : courses de fond, trois voies d'inscription}
\textbf{Énoncé :} Sur Internet, inscription progressive ; sur place, tout est contrôlé avant la course ; par courrier, un document peut manquer ; pour triathlon/vélo, c'est le club qui inscrit et la licence suffit.\\
\textbf{Solution raisonnée :} Une classe \texttt{Inscription} a un statut cohérent entre classes et états ; les variantes du parcours peuvent justifier des \textbf{gardes} selon le mode et la discipline. Le diagramme d'états n'est pas trois lignes de succès ; il a des rejets, des retours vers l'étape attendue et des délais distincts. En séquence N1, utiliser \texttt{alt} pour les voies pertinentes et les exceptions ; en N3 ne jamais laisser le formulaire contrôler seul l'âge ou la validité du certificat.
\source{Études de cas 1/2/3 ; Organisation de courses de fond.}

\section{Fiches express : juste avant l'épreuve}
\begin{tabularx}{\linewidth}{@{}p{35mm}X@{}}
\toprule
Si l'énoncé dit\ldots & Je vérifie\ldots \\
\midrule
« Le client peut\ldots » & Objectif métier / acteur / frontière \\
« Si\ldots sinon\ldots » & alt ou garde d'état ; pas une ligne oubliée \\
« Pour chaque\ldots » & loop en séquence / multiplicité adaptée \\
« Après 15 jours\ldots » & after(durée), garde et statut d'expiration \\
« Un ou plusieurs\ldots » & 1..* du côté des multiples \\
« Au plus un\ldots » & 0..1 (ou 1 si obligatoire) \\
« Au même instant\ldots » & contrainte temporelle / unicité, pas seulement cardinalité \\
« Le partenaire transmet\ldots » & acteur externe, message N1 \\
« écran / bouton\ldots » & dialogue DCP / vue N3, pas classe du domaine \\
« tarif selon\ldots » & règle/calcul, souvent classe Tarif \\
« en cas de refus\ldots » & exception N1 / N3 / états / postcondition \\
\bottomrule
\end{tabularx}

\begin{memo}
\textbf{Dernier contrôle en 15 minutes :} (1) frontière/acteurs ; (2) \texttt{include} et \texttt{extend} ; (3) cardinalités aux deux bouts ; (4) aucune classe d'interface dans le domaine ; (5) étapes alternatives ; (6) délais et retours en états ; (7) DCP sans lien dialogue--entité ; (8) N1/N3 synchronisés ; (9) règles importantes et hypothèses écrites ; (10) schémas lisibles et légendés.
\end{memo}

\clearpage% Corrigé pédagogique original, fondé sur « Organisation de courses de fond » (6 p.).
% Aucun diagramme du professeur reproduit ; choix de modélisation justifiés.
\section{Cas final : organisation de courses de fond -- correction complète des 9 travaux}
\begin{memo}
Cette partie est la \textbf{référence de travail prioritaire}. Elle répond \textbf{aux neuf demandes} du dernier TD. Le texte du professeur fait autorité. Ce corrigé est \textbf{une proposition argumentée}, pas une correction officielle. Il faut déclarer nos hypothèses H1--H4 quand les règles ne sont pas explicites.
\end{memo}
\textbf{Faits à retenir :} course libre dès 16 ans, semi-marathon dès 18 ans, marathon dès 20 ans ; triathlon et vélo : inscription par club, en ligne ou par courrier, avec licence au lieu de certificat. Préinscription internet jusqu'à J--15, sous réserve de places. Sur place : dossier terminé au moins 1 h avant le départ, pas de triathlon/vélo. Par courrier : terminé au moins J--5. Paiement invalide ou certificat non conforme : annulation en ligne ; par courrier, demande de document et maintien à l'étape adaptée. Club classé par ses cinq meilleurs temps, sinon disqualifié (au plus quatre arrivants). SMS en option payante à 1 euro. Transpondeur : numéro de dossard, temps et passages, restitution obligatoire.

\begin{warningbox}
\textbf{Hypothèses déclarées :} H1 : un agent du secrétariat saisit les inscriptions reçues sur papier ou au guichet ; H2 : on sépare l'état de l'\texttt{Inscription} du support de paiement ; H3 : le statut d'inscription suit des étapes communes, sauf les parcours plus rapides qui sautent plusieurs étapes ; H4 : la catégorie d'âge et l'admissibilité sont vérifiées par \texttt{Epreuve} (le texte n'impose aucune classe responsable). Ne transformez pas ces choix en « faits du sujet ».
\end{warningbox}

\subsection{1. Cas d'utilisation : acteurs, objectifs et paquetages}
\textbf{Acteurs primaires :} Concurrent, Club (pour vélo et triathlon), AgentSecrétariat, ServiceCourses, Administrateur (si gestion technique précisée). \textbf{Acteurs secondaires :} Banque/PrestataireCarte, SystèmeChronométrage (antennes et puces externes, sous H5 si le périmètre du logiciel ne les inclut pas), ServiceSMS (si externe). Un transpondeur passif n'est \textbf{pas automatiquement} acteur ; seul le dispositif qui échange effectivement des messages avec le système peut l'être.
\begin{center}
\begin{tikzpicture}[scale=.84,transform shape,node distance=9mm and 17mm]
\node[actor] (cc) at (0,2.8) {Concurrent};\node[actor] (club) at (0,.8) {Club};
\node[actor] (sec) at (0,-1.3) {Secrétariat};\node[actor] (srv) at (0,-3.3) {Service courses};
\node[usecase] (online) at (4.5,3.2) {S'inscrire sur internet};
\node[usecase] (cour) at (4.5,1.7) {Inscrire par courrier};
\node[usecase] (place) at (4.5,.2) {Inscrire sur place};
\node[usecase] (consult) at (4.5,-1.3) {Rechercher un classement};
\node[usecase] (classer) at (4.5,-2.8) {Enregistrer classement};
\node[usecase] (result) at (4.5,-4.3) {Enregistrer les temps};
\node[actor] (bank) at (8.7,3.4) {Banque};
\node[actor] (chrono) at (8.7,-4.3) {Système de\\chronométrage};
\draw[navy] (2.7,-5.05) rectangle (7.1,4.05);
\node[font=\scriptsize,fill=white] at (4.9,4.04){Système d'organisation des courses};
\draw[navy] (cc)--(online);\draw[navy](cc)--(consult);
\draw[navy](club)--(online);\draw[navy](club)--(cour);
\draw[navy](sec)--(cour);\draw[navy](sec)--(place);
\draw[navy](srv)--(classer);\draw[navy](srv)--(result);
\draw[navy](bank)--(online);\draw[navy](chrono)--(result);
\end{tikzpicture}
\end{center}
\textbf{Paquetages fonctionnels :} \texttt{Inscriptions}, \texttt{EpreuvesEtClubs}, \texttt{Classements}, \texttt{Chronometrage}. Un paquetage de cas n'est pas une étape du scénario. Les cas « consulter conditions » et « télécharger formulaire » appartiennent à l'espace public ; « gérer récompenses » au ServiceCourses. \textbf{Attention :} un Club n'est pas un « Concurrent amélioré » ; c'est un rôle institutionnel différent.

\subsection{2. Description textuelle détaillée : inscription sur internet}
\begin{memo}
\textbf{Méthode exacte d'ASI\_C1, page 36 :} Identification (nom, but, acteurs) ; Séquencement (déclencheur, préconditions, enchaînement nominal, postconditions) ; variantes et exceptions ; éventuellement exigences non fonctionnelles et IHM. Le TD de course de fond demande explicitement ce travail : ne pas le réduire à un diagramme.
\end{memo}
\begin{tabularx}{\linewidth}{@{}p{40mm}X@{}}
\toprule
\textbf{Nom du cas} & S'inscrire à une course sur internet.\\
\textbf{But} & Obtenir une inscription définitive après paiement accepté et vérification des pièces.\\
\textbf{Acteur principal} & Concurrent pour une course individuelle ; Club de la discipline pour triathlon ou vélo.\\
\textbf{Acteurs secondaires} & Banque (carte) ; Secrétariat (chèque, certificat ou licence).\\
\textbf{Déclencheur} & Le concurrent choisit une course et demande son inscription depuis le site web.\\
\textbf{Préconditions} & Une épreuve est publiée et accessible via le site. On \emph{contrôle pendant le scénario} la date, les places et l'admissibilité.\\
\textbf{Postcondition succès} & Une \texttt{Inscription} est \texttt{DEFINITIVE} ; paiement accepté et justificatif conforme.\\
\textbf{Postconditions sans succès} & Avant la préinscription, aucun dossier si quota/date/âge refusé. Après paiement invalide ou pièce non conforme : \texttt{ANNULEE}, concurrent prévenu. Dossier temporairement \texttt{PREINSCRIPTION} ou \texttt{FINANCIERE} en attendant les pièces.\\
\bottomrule
\end{tabularx}
\textbf{Enchaînement nominal (course individuelle ordinaire, carte) :}
\begin{enumerate}[label=\textbf{\arabic*.}]
\item Le concurrent sélectionne la course. Le système affiche ses conditions, sa date et son tarif.
\item Le système vérifie le quota, la limite J--15 et la condition d'âge liée à la discipline.
\item Le concurrent fournit ou confirme ses données personnelles. Le système retrouve les données déjà connues, si elles existent.
\item Le système crée l'inscription au statut \texttt{PREINSCRIPTION}.
\item Le concurrent choisit de payer par carte. Le système transmet la demande de paiement à la banque.
\item La banque accepte. Le système fait passer l'inscription au statut \texttt{FINANCIERE}.
\item Le concurrent envoie par courrier un certificat médical de moins d'un an. Le secrétariat le reçoit et le contrôle.
\item Si le certificat est conforme, le système passe à \texttt{DEFINITIVE} et le concurrent est informé.
\end{enumerate}
\textbf{Enchaînements alternatifs :}
\begin{itemize}
\item \textbf{A1 [étape 3 : concurrent déjà connu] :} le système propose les anciennes informations ; le concurrent les valide ou les corrige ; retour à l'étape 4.
\item \textbf{A2 [étapes 5--6 : chèque] :} le concurrent envoie son paiement par courrier ; le secrétariat le contrôle ; le statut devient \texttt{FINANCIERE} quand le chèque est accepté. Le certificat est vérifié ensuite.
\item \textbf{A3 [triathlon / vélo] :} le \texttt{Club} inscrit son membre par internet ; on contrôle l'adhésion au club spécialisé ; la \texttt{Licence} conforme remplace le certificat médical ; il n'y a \textbf{pas} d'inscription sur place.
\end{itemize}
\textbf{Exceptions :}
\begin{itemize}
\item \textbf{E1 [étape 2 : places insuffisantes ou date J--15 dépassée] :} refus avant préinscription. La condition d'âge empêche également l'inscription aux courses individuelles non autorisées.
\item \textbf{E2 [étape 6 : paiement non valide] :} inscription annulée et concurrent informé.
\item \textbf{E3 [étape 8 : certificat non conforme] :} inscription annulée et concurrent informé.
\end{itemize}
\textbf{Règles métier et remarques :} libre \(\geq16\) ans, semi \(\geq18\), marathon \(\geq20\) ; quota et date J--15 ; certificat médical datant de moins d'un an (ou licence du club spécialisé). \textbf{N'inventez pas} un upload internet de certificat ni une double autorisation/débit : ici la validation définitive est liée à la vérification du justificatif.\par
\textbf{Contrainte non fonctionnelle possible (suggestion, non donnée mot pour mot) :} limiter l'accès aux données personnelles et médicales aux personnes autorisées.
\begin{warningbox}
Un \textbf{paiement accepté} ne signifie pas \textbf{inscription définitive}. La préinscription, le paiement et le contrôle du justificatif sont \emph{trois étapes différentes}. Un courrier incomplet, lui, n'est pas automatiquement annulé : le secrétariat relance le concurrent et conserve l'étape en cours.
\end{warningbox}
\begin{tabularx}{\linewidth}{@{}p{29mm}XXX@{}}
\toprule
Pour comparer & Internet & Sur place & Courrier \\
\midrule
Date clé & Préinscription au plus tard J--15 & Dossier complet 1 h avant départ & Dossier complet au plus tard J--5\\
Qui intervient ? & Concurrent (ou club), banque, secrétariat & Concurrent, secrétariat & Concurrent/club, secrétariat\\
Mode paiement & Carte ou chèque & Espèces, carte ou chèque & Règlement joint, souvent chèque\\
Justificatif & Certificat envoyé par courrier ; licence si club & Certificat remis et contrôlé & Certificat joint ; licence si club\\
Si pièce manque & Dossier attend tant que possible & Validation immédiate impossible & Conserver étape et demander pièce\\
Triathlon/vélo & Club uniquement & \textbf{Non admis} & Club uniquement\\
\bottomrule
\end{tabularx}
\source{Organisation de courses de fond, pages 3--4 ; ASI\_C1, page 36.}

\subsection{3. Les trois séquences système N1 d'inscription}
\textbf{Règle :} un seul participant \texttt{:SystèmeCourses}, plus acteurs externes. Aucun contrôleur, écran ou DAO. Chaque case \texttt{alt} correspond à un véritable embranchement.
\begin{center}
\begin{tikzpicture}[x=.94cm,y=.49cm,scale=.91,transform shape]
\foreach \x/\t in {0/Concurrent,3.5/SystèmeCourses,7.3/Banque,11/Secrétariat}{\node[dbox] at (\x,0){\t};\draw[life](\x,-.5)--(\x,-20.8);}
\draw[assoc](0,-1)--node[above,font=\scriptsize]{1: choisirEpreuve()}(3.5,-1);
\draw[<-,dashed](0,-2)--node[above,font=\scriptsize]{2: conditions, quota, tarif}(3.5,-2);
\draw[teal](-.6,-2.6)rectangle(4.3,-4.6);
\node[font=\tiny,anchor=west] at (-.4,-2.87){break [quota/date/âge non valides]};
\draw[<-,dashed](0,-3.73)--node[above,font=\tiny]{3: refus motivé, pas de dossier}(3.5,-3.73);
\draw[assoc](0,-5.2)--node[above,font=\scriptsize]{4: saisirOuConfirmerDonnees()}(3.5,-5.2);
\draw[<-,dashed](0,-6.2)--node[above,font=\scriptsize]{5: préinscription enregistrée}(3.5,-6.2);
\draw[assoc](0,-7.2)--node[above,font=\scriptsize]{6: choisirPaiement()}(3.5,-7.2);
\draw[teal](-.6,-7.7)rectangle(11.9,-12.25);
\node[font=\tiny,anchor=west] at (-.4,-7.94){alt [carte]};
\draw[assoc](3.5,-8.75)--node[above,font=\scriptsize]{7a: demanderPaiement()}(7.3,-8.75);
\draw[<-,dashed](3.5,-9.7)--node[above,font=\scriptsize]{8a: accepté/refusé}(7.3,-9.7);
\draw[teal,densely dashed](-.6,-10.18)--(11.9,-10.18);
\node[font=\tiny,anchor=west] at (-.4,-10.4){[chèque]};
\draw[assoc](11,-11.48)--node[above,font=\scriptsize]{7b: saisirRésultatContrôleChèque()}(3.5,-11.48);
\draw[teal](-.6,-12.75)rectangle(11.9,-16.2);
\node[font=\tiny,anchor=west] at (-.4,-13){alt [paiement invalide]};
\draw[<-,dashed](0,-13.85)--node[above,font=\tiny]{9a: ANNULEE, avertissement}(3.5,-13.85);
\draw[teal,densely dashed](-.6,-14.35)--(11.9,-14.35);
\node[font=\tiny,anchor=west] at (-.4,-14.58){[paiement accepté]};
\draw[<-,dashed](0,-15.4)--node[above,font=\tiny]{9b: FINANCIERE, attendre pièce}(3.5,-15.4);
\draw[assoc](11,-16.75)--node[above,font=\scriptsize]{10: contrôlerCertificatOuLicence()}(3.5,-16.75);
\draw[teal](-.6,-17.2)rectangle(11.9,-20.45);
\node[font=\tiny,anchor=west] at (-.4,-17.43){alt [justificatif conforme]};
\draw[<-,dashed](0,-18.25)--node[above,font=\tiny]{11a: DEFINITIVE, avertissement}(3.5,-18.25);
\draw[teal,densely dashed](-.6,-18.66)--(11.9,-18.66);
\node[font=\tiny,anchor=west] at (-.4,-18.87){[justificatif non conforme]};
\draw[<-,dashed](0,-19.76)--node[above,font=\tiny]{11b: ANNULEE, avertissement}(3.5,-19.76);
\end{tikzpicture}
\end{center}
\textbf{N1 internet -- lecture exacte :} les deux \texttt{alt} « paiement » et « justificatif » sont deux décisions \textbf{successives}. La seconde n'a lieu que si le paiement a été accepté. Le secrétariat peut saisir l'information après réception du certificat \emph{par courrier}. Le club remplace l'acteur principal et présente une \emph{licence} pour triathlon/vélo. Les réponses aux contrôles font partie du N1, même si les échanges postaux eux-mêmes ont lieu hors système.
\begin{center}\begin{tikzpicture}[x=1cm,y=.72cm]
\foreach \x/\t in {0/Concurrent,5/SystèmeCourses,10/Secrétariat}{\node[dbox] at (\x,0){\t};\draw[life](\x,-.5)--(\x,-8.6);}
\draw[assoc](0,-1)--node[above,font=\scriptsize]{1: remettreFormulaireEtPièces()}(10,-1);
\draw[assoc](10,-2)--node[above,font=\scriptsize]{2: saisirInscription()}(5,-2);
\draw[densely dashed,-Stealth](5,-3)--node[above,font=\scriptsize]{3: résultat contrôles}(10,-3);
\draw[teal] (4,-3.6) rectangle (10.75,-8.2);
\node[font=\scriptsize,anchor=west] at (4.2,-3.95){alt [valide, avant départ moins 1h et places]};
\draw[assoc](10,-4.7)--node[above,font=\scriptsize]{4a: enregistrer définitivement()}(5,-4.7);
\draw[densely dashed,-Stealth](5,-5.7)--node[above,font=\scriptsize]{5a: confirmer inscription}(10,-5.7);
\draw[teal,densely dashed](4,-6.1)--(10.75,-6.1);
\node[font=\scriptsize,anchor=west] at (4.2,-6.4){[sinon]};
\draw[densely dashed,-Stealth](5,-7.25)--node[above,font=\scriptsize]{4b: rejet et motif}(10,-7.25);
\end{tikzpicture}\end{center}
\textbf{N1 sur place :} la branche est refusée aussi pour triathlon/vélo. Pour un paiement carte, la Banque reste un acteur secondaire distinct (ajouter son échange avant confirmation).
\begin{center}\begin{tikzpicture}[x=1cm,y=.72cm]
\foreach \x/\t in {0/Concurrent,5/Secrétariat,10/SystèmeCourses}{\node[dbox] at (\x,0){\t};\draw[life](\x,-.5)--(\x,-8.1);}
\draw[assoc](0,-1)--node[above,font=\scriptsize]{1: envoyerDossierPapier()}(5,-1);
\draw[assoc](5,-2)--node[above,font=\scriptsize]{2: saisirEtContrôlerDossier()}(10,-2);
\draw[densely dashed,-Stealth](10,-3)--node[above,font=\scriptsize]{3: état du dossier}(5,-3);
\draw[teal](4.2,-3.6)rectangle(10.6,-7.8);
\node[font=\scriptsize,anchor=west] at (4.4,-3.9){alt [complet, reçu avant J--5]};
\draw[densely dashed,-Stealth](10,-4.7)--node[above,font=\scriptsize]{4a: inscription définitive}(5,-4.7);
\draw[teal,densely dashed](4.2,-5.15)--(10.6,-5.15);
\node[font=\scriptsize,anchor=west] at (4.4,-5.46){[pièce manquante ou paiement non conforme]};
\draw[densely dashed,-Stealth](10,-6.25)--node[above,font=\scriptsize]{4b: maintien état adéquat}(5,-6.25);
\draw[assoc](5,-7.05)--node[above,font=\scriptsize]{5b: demander pièce au concurrent}(0,-7.05);
\end{tikzpicture}\end{center}
\textbf{N1 courrier :} la règle J--5 s'applique à la validation finale ; une pièce manquante peut être relancée, sans annulation automatique immédiate. Pour une inscription de club, le courrier vient du Club.

\subsection{4. Classes du domaine : noms, propriétés, opérations, relations}
\begin{center}\begin{tikzpicture}[scale=.86,transform shape]
\node[dbox,align=left,minimum width=35mm] (co) at (0,3) {\textbf{Concurrent}\\nom : String\\naissance : Date\\sexe : Sexe\\estMajeurPour(e) : Bool};
\node[dbox,align=left,minimum width=37mm] (ins) at (5,3) {\textbf{Inscription}\\date : Date\\dossard : String\\statut : StatutInscription\\valider() / annuler()};
\node[dbox,align=left,minimum width=40mm] (ep) at (10,3) {\textbf{Epreuve}\\dateDépart : DateTime\\quotaMax : Integer\\tarif : Money\\estAdmissible(c) : Bool};
\node[dbox,align=left,minimum width=35mm] (club) at (0,0) {\textbf{Club}\\nom : String\\discipline : Discipline};
\node[dbox,align=left,minimum width=35mm] (pa) at (5,0) {\textbf{Paiement}\\montant : Money\\mode : ModePaiement\\estAccepté() : Bool};
\node[dbox,align=left,minimum width=37mm] (just) at (10,0) {\textbf{Justificatif}\\dateValidité : Date\\type : TypeJustif\\estConforme() : Bool};
\node[dbox,align=left,minimum width=35mm] (temps) at (0,-3) {\textbf{Passage}\\instant : DateTime\\tempsNet : Durée};
\node[dbox,align=left,minimum width=35mm] (trans) at (5,-3) {\textbf{Transpondeur}\\codeDossard : String\\estRestitué : Bool};
\node[dbox,align=left,minimum width=37mm] (cla) at (10,-3) {\textbf{Classement}\\position : Integer\\tempsFinal : Durée\\calculerPosition()};
\draw[navy](co)--node[above,font=\scriptsize,pos=.10]{1}node[above,font=\scriptsize,pos=.86]{0..*}(ins);
\draw[navy](ins)--node[above,font=\scriptsize,pos=.10]{0..*}node[above,font=\scriptsize,pos=.86]{1}(ep);
\draw[navy](co)--node[left,font=\scriptsize]{0..1 / 0..*}(club);
\draw[navy](ins)--node[left,font=\scriptsize]{1 / 0..1}(pa);
\draw[navy](ins) to[bend left=13] node[right,font=\scriptsize,pos=.5]{1 / 0..1} (just);
\draw[navy](ins)--node[above,sloped,font=\scriptsize]{0..1}(trans);
\draw[navy](trans)--node[above,font=\scriptsize]{0..*}(temps);
\draw[navy](ins)--node[above,sloped,font=\scriptsize]{0..1}(cla);
\end{tikzpicture}\end{center}
\textbf{Précision critique :} les multiplicités abrégées sur les liens obliques ci-dessus se lisent comme des \textbf{hypothèses de modélisation}, à détailler et justifier : une Inscription concerne exactement un Concurrent et une Epreuve ; un Concurrent peut avoir plusieurs Inscription ; un Paiement est rattaché au plus à une inscription individuelle dans ce modèle simplifié. Une catégorie de classement est une \textbf{valeur métier}, pas une nouvelle personne. Un Transpondeur peut être réaffecté dans le temps : les passages restent attachés à l'Inscription (ajouter une association précise en version de production).
\textbf{Classes supplémentaires utiles :} \texttt{Edition}, \texttt{Troncon}, \texttt{Antenne}, \texttt{Recompense}, \texttt{Notification}, \texttt{ClubClassement}. Ne pas stocker un rang figé avant l'arrivée : la position est calculée depuis les temps. \textbf{Classe d'association} : \texttt{Inscription} porte les attributs du lien Concurrent--Epreuve. Les contraintes d'âge et de quotas appartiennent à \texttt{Epreuve} ; les statuts à \texttt{Inscription}.

\textbf{Multiplicités à maîtriser pour ne pas perdre de points :}\par\smallskip
\begin{tabularx}{\linewidth}{@{}p{55mm}p{20mm}X@{}}
\toprule
Association (de A vers B) & A vers B & Lecture et hypothèse \\
\midrule
\texttt{Concurrent -- Inscription} & \texttt{0..* / 1} & Un concurrent peut s'inscrire plusieurs fois ; une inscription correspond à un concurrent.\\
\texttt{Epreuve -- Inscription} & \texttt{0..* / 1} & Une épreuve reçoit plusieurs inscriptions ; chacune vise une seule épreuve.\\
\texttt{Inscription -- Paiement} & \texttt{0..1 / 1} & Un dossier peut attendre le règlement ; un règlement enregistré appartient à un dossier (H2).\\
\texttt{Inscription -- Justificatif} & \texttt{0..* / 1} & Un dossier peut avoir des pièces successives ; chaque pièce est reliée à un dossier (H7).\\
\texttt{Epreuve -- Troncon} & \texttt{1..* / 1} & Parcours divisé en tronçons ; un tronçon est rattaché à son épreuve (H8).\\
\texttt{Troncon -- Antenne} & \texttt{1 / 1} & Une antenne à la fin de chaque tronçon : simplification du TD (H8).\\
\texttt{Transpondeur -- Passage} & \texttt{0..* / 1} & Plusieurs passages peuvent provenir d'une puce ; un passage porte son émetteur.\\
\texttt{Inscription -- Classement} & \texttt{0..1 / 1} & Un classement individuel est lié à une inscription terminée ; aucun rang avant résultat (H9).\\
\bottomrule
\end{tabularx}
\textbf{Sens de la table :} \texttt{A -- B : 0..* / 1} signifie qu'un A peut avoir plusieurs B et qu'un B est rattaché à un seul A. Ce sont des \textbf{choix de modélisation justifiés}, pas des multiplicités toutes écrites dans l'énoncé.\\
\textbf{Ne pas confondre :} \texttt{Inscription} rattache un \texttt{Concurrent} à une \texttt{Epreuve}. Un \texttt{Club} n'est pas un type de \texttt{Concurrent}. Le classement \textbf{par club} est calculé avec les cinq meilleurs temps ; si au plus quatre membres terminent, le club est disqualifié.

\textbf{Vue métier complémentaire -- chronométrage et récompenses :}
\begin{center}
\begin{tikzpicture}[scale=.86,transform shape]
\node[dbox] (ep) at (0,1.6) {Epreuve};
\node[dbox] (tr) at (3,1.6) {Troncon};
\node[dbox] (ant) at (6,1.6) {Antenne};
\node[dbox] (pas) at (9,1.6) {Passage};
\node[dbox] (trans) at (9,-.1) {Transpondeur};
\node[dbox] (ins) at (3,-.1) {Inscription};
\node[dbox] (clas) at (3,-1.8) {Classement};
\node[dbox] (rew) at (7.5,-1.8) {Recompense};
\draw[navy](ep)--node[above,font=\tiny,pos=.25]{1}node[above,font=\tiny,pos=.78]{1..*}(tr);
\draw[navy](tr)--node[above,font=\tiny,pos=.22]{1}node[above,font=\tiny,pos=.80]{1}(ant);
\draw[navy](ant)--node[above,font=\tiny,pos=.22]{1}node[above,font=\tiny,pos=.80]{0..*}(pas);
\draw[navy](trans)--node[right,font=\tiny]{1 / 0..*}(pas);
\draw[navy](ins)--node[left,font=\tiny]{1 / 0..1}(clas);
\draw[navy](clas)--node[above,font=\tiny]{0..* (H9)}(rew);
\end{tikzpicture}
\end{center}
\textbf{Noter H8--H9 :} un transpondeur ne remplace pas la classe \texttt{Inscription}. Pour relier chaque temps au concurrent de la bonne course, une affectation de puce à inscription est nécessaire (association ou classe d'association si l'affectation porte des dates).\par
\subsection{5. Paquetages et dépendances}
\begin{center}\begin{tikzpicture}[node distance=12mm and 20mm]
\node[cbox] (s) {\textbf{Inscriptions}\\Concurrent, Inscription, Paiement\\Justificatif};
\node[cbox,right=of s] (ev) {\textbf{Epreuves et clubs}\\Epreuve, Club, Edition};
\node[cbox,below=of s] (ch) {\textbf{Chronométrage}\\Transpondeur, Passage\\Troncon, Antenne};
\node[cbox,right=of ch] (res) {\textbf{Classements}\\Classement, Récompense};
\draw[umlarrow](s)--(ev);\draw[umlarrow](res)--(s);\draw[umlarrow](res)--(ch);\draw[umlarrow](ch)--(ev);
\end{tikzpicture}\end{center}
\textbf{Critère :} découpage par responsabilité métier ; les dépendances pointent vers le paquetage \emph{utilisé}. Pas de dépendance circulaire. Si une opération impose un cycle, déplacer un concept dans un noyau partagé ou revoir le sens du lien.

\subsection{6. États de l'inscription individuelle avec déclencheurs complets}
\begin{center}\begin{tikzpicture}[node distance=15mm and 20mm]
\node[state] (pre) {Préinscription};
\node[state,right=of pre] (fin) {Inscription financière};
\node[state,right=of fin] (def) {Inscription définitive};
\node[state,below=of fin] (ann) {Annulée};
\node[circle,fill=navy,inner sep=2pt,left=10mm of pre] (ini) {};
\node[draw,circle,minimum size=7mm,inner sep=0pt,right=14mm of def] (end) {}; \node[fill=navy,circle,minimum size=3mm,inner sep=0pt] at (end.center) {};
\draw[assoc](ini)--node[above,font=\tiny]{préinscrire [quota et date OK]}(pre);
\draw[assoc](pre)--node[above,font=\tiny]{paiementAccepté / passerFinancière()}(fin);
\draw[assoc](fin)--node[above,font=\tiny]{pièceConforme / confirmer()}(def);
\draw[assoc](def)--(end);
\draw[assoc](pre) to[bend right=28] node[below,font=\tiny]{paiementRefusé / annuler()} (ann);
\draw[assoc](fin)--node[right,font=\tiny]{certificatRefusé / annuler()}(ann);
\end{tikzpicture}\end{center}
\textbf{Énumération :} \texttt{StatutInscription = \{PREINSCRIPTION, FINANCIERE, DEFINITIVE, ANNULEE\}}. Dans le parcours \textbf{courrier}, si une pièce manque on conserve PREINSCRIPTION ou FINANCIERE et on peut reprendre le contrôle : représenter une \textbf{boucle de régularisation} \texttt{pièceReçue [conforme] / continuer()}. Dans le parcours \textbf{sur place}, trois étapes peuvent être franchies dans la même opération (pas forcément trois écrans). \textbf{Délai :} J--15 bloque l'entrée internet, J--5 la validation courrier et 1 h le guichet ; ne pas inventer un \texttt{after(24h)} inexistant dans ce TD. Le mot \texttt{after(durée)} du cours est pertinent seulement quand un délai relatif déclenche réellement un changement d'état.

\subsection{7. Deux DCP : inscription internet et recherche classement}
\begin{center}\begin{tikzpicture}[node distance=10mm and 12mm]
\node[actor] (a) {Concurrent};
\node[cbox,right=of a] (v) {«dialogue»\\PageInscription};
\node[cbox,right=of v] (c) {«contrôle»\\ControleInscription};
\node[cbox,right=of c,yshift=26mm] (e) {«entité»\\Epreuve};
\node[cbox,right=of c,yshift=13mm] (i) {«entité»\\Inscription};
\node[cbox,right=of c,yshift=0mm] (con) {«entité»\\Concurrent};
\node[cbox,right=of c,yshift=-13mm] (p) {«entité»\\Paiement};
\node[cbox,right=of c,yshift=-26mm] (j) {«entité»\\Justificatif};
\draw[navy](a)--(v);\draw[navy](v)--(c);\draw[navy](c)--(e);\draw[navy](c)--(i);\draw[navy](c)--(con);\draw[navy](c)--(p);\draw[navy](c)--(j);
\end{tikzpicture}\end{center}
\textbf{DCP inscription :} cinq entités du modèle : \texttt{Concurrent}, \texttt{Epreuve}, \texttt{Inscription}, \texttt{Paiement}, \texttt{Justificatif}. \texttt{PagePaiement} et \texttt{PageRecapitulatif} sont des dialogues \textbf{supposés} (H6), non fournis comme maquettes. La seule maquette réellement donnée est celle de recherche des classements. Le contrôle orchestre ; les entités valident la règle métier.
\begin{center}\begin{tikzpicture}[node distance=9mm and 12mm]
\node[actor] (a) {Visiteur};\node[cbox,right=of a] (v) {«dialogue»\\PageRecherche};
\node[cbox,right=of v] (c) {«contrôle»\\ControleRecherche};
\node[cbox,right=of c,yshift=16mm] (e) {«entité»\\Concurrent};
\node[cbox,right=of c] (i) {«entité»\\Inscription};
\node[cbox,right=of c,yshift=-16mm] (cl) {«entité»\\Classement};
\draw[navy](a)--(v);\draw[navy](v)--(c);\draw[navy](c)--(e);\draw[navy](c)--(i);\draw[navy](c)--(cl);
\end{tikzpicture}\end{center}
\textbf{DCP classement :} les critères exacts de la maquette sont \texttt{Dossard, Nom, Prénom, Sexe, Catégorie}. Les entités doivent reprendre exactement les noms du domaine : \texttt{Concurrent, Inscription, Classement}. \textbf{Interdit :} un lien direct Dialogue--Entité. Un \texttt{ControleRecherche} distinct est justifié car c'est un autre cas d'utilisation.

\subsection{8. Communication pour rechercher un classement}
\begin{center}\begin{tikzpicture}[node distance=15mm and 22mm]
\node[actor] (v) {Visiteur};\node[cbox,right=of v] (vue) {:PageRecherche};
\node[cbox,right=of vue] (ctrl) {:ControleRecherche};
\node[cbox,right=of ctrl] (dom) {:Classement};
\draw[navy](v)--node[above,font=\tiny]{1: saisirCritères()}(vue);
\draw[navy](vue)--node[above,font=\tiny]{2: rechercher()}node[below,font=\tiny]{5: afficher()}(ctrl);
\draw[navy](ctrl)--node[above,font=\tiny]{3: obtenirRésultats()}node[below,font=\tiny]{4: résultats}(dom);
\end{tikzpicture}\end{center}
\textbf{Complément de modèle :} \texttt{ControleRecherche} peut demander au \texttt{Concurrent} les identités et à \texttt{Inscription} les dossards. L'ordre des messages est écrit \textbf{sur les liens}, pas sur un axe du temps. Une numérotation détaillée des appels imbriqués \texttt{1, 1.1, 1.2} est aussi valide si cohérente.

\subsection{9. Séquence N3 MVC : inscription sur place}
\begin{center}\begin{tikzpicture}[x=.82cm,y=.54cm]
\foreach \x/\t in {0/Agent,3/PageGuichet,6/ControleInscription,9/Epreuve,12/Inscription}{\node[dbox,minimum width=21mm] at (\x,0){\t};\draw[life](\x,-.5)--(\x,-12.5);}
\draw[assoc](0,-1)--node[above,font=\tiny]{1: saisirDossier()}(3,-1);
\draw[assoc](3,-2)--node[above,font=\tiny]{2: inscrireSurPlace()}(6,-2);
\draw[assoc](6,-3)--node[above,font=\tiny]{3: verifierQuotaAgeDélai()}(9,-3);
\draw[densely dashed,-Stealth](9,-4)--node[above,font=\tiny]{4: résultat}(6,-4);
\draw[teal](5.25,-4.6)rectangle(12.8,-11.9);
\node[font=\tiny,anchor=west] at (5.45,-4.9){alt [admissible]};
\draw[assoc](6,-5.8)--node[above,font=\tiny]{5a: contrôlerPaiementPièce()}(12,-5.8);
\draw[densely dashed,-Stealth](12,-6.75)--node[above,font=\tiny]{6a: statut contrôle}(6,-6.75);
\draw[assoc](6,-7.7)--node[above,font=\tiny]{7a: confirmer()}(12,-7.7);
\draw[assoc](6,-8.6)--node[above,font=\tiny]{8a: afficherConfirmation()}(3,-8.6);
\draw[teal,densely dashed](5.25,-9.2)--(12.8,-9.2);
\node[font=\tiny,anchor=west] at (5.45,-9.5){[quota atteint / délai / justificatif invalide]};
\draw[assoc](6,-10.4)--node[above,font=\tiny]{5b: afficherRefus(motif)}(3,-10.4);
\end{tikzpicture}\end{center}
\textbf{Participants métier à conserver dans une version exhaustive :} \texttt{Concurrent}, \texttt{Paiement}, \texttt{Justificatif}, au besoin \texttt{Banque} pour carte et \texttt{Club} en exclusion du guichet. Le contrôle appelle les \emph{opérations métier des entités}; il ne calcule pas lui-même les seuils d'âge et de délai. L'interface ne parle jamais directement à \texttt{Epreuve} ni à \texttt{Inscription}. Le N3 reprend les interactions externes du N1.

\textbf{N3 détaillée -- messages qui doivent exister, même si le schéma précédent est compact :}
\begin{longtable}{@{}p{13mm}p{52mm}p{80mm}@{}}
\toprule
No & Émetteur vers destinataire & Message et rôle \\
\midrule\endhead
1 & Agent \(\to\) \texttt{PageGuichet} & \texttt{saisirDossier(...)}. Vue : saisie des informations et pièces.\\
2 & Vue \(\to\) \texttt{ControleInscription} & \texttt{inscrireSurPlace(...)} ; le contrôle orchestre.\\
3 & Contrôle \(\to\) \texttt{Concurrent} & \texttt{retrouverOuCreer(...)} ; réutiliser les données d'une précédente course.\\
4 & Contrôle \(\to\) \texttt{Epreuve} & \texttt{verifierAdmissibilite()}. Règles du métier : \textbf{1 h}, quota et âge ; \textbf{pas triathlon/vélo}.\\
5a & \texttt{Epreuve} \(\to\) Contrôle & \texttt{refus(motif)} ; contrôle \(\to\) vue \texttt{afficherRefus(motif)}. Fin du cas si inadmissible.\\
5b & Contrôle \(\to\) \texttt{Paiement} & \texttt{verifier(mode,montant)} ; en cas de carte, solliciter aussi la \texttt{Banque} externe.\\
6b & Contrôle \(\to\) \texttt{Justificatif} & \texttt{estConforme()} pour le certificat médical récent.\\
7b & Contrôle \(\to\) \texttt{Inscription} & \texttt{enregistrerEtValider()} seulement si \textbf{tous les contrôles sont OK}.\\
8b & Contrôle \(\to\) \texttt{PageGuichet} & \texttt{afficherConfirmation(dossard)} puis vue \(\to\) Agent.\\
\bottomrule
\end{longtable}
\textbf{Branches à écrire dans la N3 :} \texttt{alt [épreuve non admissible] / [admissible]} ; puis \texttt{alt [paiement invalide OU certificat invalide] / [tous conformes]}. En cas d'erreur, Le contrôle informe la vue ; le dossier ne devient pas \texttt{DEFINITIVE}. \textbf{H10 :} la manière de tracer une saisie incomplète n'est pas spécifiée dans le texte ; ne pas inventer une notification SMS pour le guichet.\par
\textbf{Cohérence forte :} conserver les mêmes noms de dialogues, contrôles et entités entre le DCP et la N3. Les noms \texttt{Concurrent}, \texttt{Epreuve}, \texttt{Paiement}, \texttt{Justificatif} et \texttt{Inscription} reprennent ceux du domaine.

\subsection{Dernière vérification des neuf réponses}
\begin{tabularx}{\linewidth}{@{}p{24mm}X@{}}\toprule
Questions & Critère de cohérence à vérifier avant de rendre \\\midrule
1--2 & Même acteur principal, même scénario, différences club / individuel claires.\\
3 & Trois canaux modélisés, cas d'erreur et retours, système noir.\\
4--5 & Noms métier et multiplicités ; aucune classe de DAO dans le domaine, pas de cycle.\\
6 & Énumération alignée et branches d'annulation/régularisation présentes.\\
7--8 & Dialogue--Contrôle--Entités, acteurs à l'extérieur, mêmes noms, messages numérotés.\\
9 & N3 MVC détaillée, mêmes événements que N1, contrôle court, règles dans les entités.\\\bottomrule\end{tabularx}\clearpage\section*{Sources et portée de cette révision}
\addcontentsline{toc}{section}{Sources et portée}
Base du document : \textit{ASI\_C1 Introduction et modèle des cas d'utilisation} ; \textit{ASI\_C2 Classes, paquetages et états} ; \textit{ASI\_C3 De l'analyse à la conception} ; \textit{ASI\_C4 Conduite de projet et agilité} ; \textit{ASI\_C5 Synthèse, IA et projet} ; \textit{Document de vision} ; \textit{Exercices et études de cas (1, 2, 3)} ; \textit{Location de véhicules} ; \textit{Organisation de courses de fond} ; \textit{SAFe Explained FR} (complément).\\
Les schémas et mini-cas de cette révision sont des \textbf{illustrations originales}, pas des corrigés officiels. Toute hypothèse non fournie par un énoncé est signalée comme telle. Le contenu couvre le noyau de la matière et les thèmes complémentaires visibles dans les ressources fournies ; il ne peut pas garantir mot pour mot le contenu d'un futur sujet.
\end{document}